\chapter{Electrostatics and Capacitance}
\label{ch:electrostatics}

\chapterorient{This chapter opens Part~V, and it is the one the rest of the part
is written against. Every modality chapter to come is a variation on a single
theme, and here we play that theme in full: we take the electrostatic
analysis---the simplest of the six---and trace it end to end, from a
configuration of conductors to the capacitance matrix that comes out, with no box
left closed. We will watch the skeleton of \cref{ch:situating} become an actual
computation: assemble the stiffness operator \emph{once}, map a family of
independent solves over it (the fixed-operator corner), and Gram-reduce the solved
fields into the answer. By the end you will have seen a complete analysis run, and
you will be able to read every later modality as \emph{this pipeline, with one
stage swapped}.}

\section{The physical question}
\label{sec:es-question}

An engineer hands us a device---a few metal conductors embedded in an insulating
dielectric---and asks a question as old as the Leyden jar: \emph{how much charge
does it store per volt?} For a single isolated conductor the answer is one number,
its capacitance. For $n$ conductors the answer is a \emph{matrix}, because the
charge on each conductor depends on the voltage of \emph{every} conductor, not
just its own. That $n\times n$ table is the \emph{capacitance matrix}
$\mat{C}$, and computing it is the entire job of the electrostatic analysis.

\begin{physicsbox}
Several conductors sit in a dielectric medium (\cref{fig:es-setup}). Each
conductor is a \emph{perfect equipotential}: a perfect conductor cannot sustain a
tangential electric field, so its whole surface sits at one voltage. We are free
to \emph{set} those voltages; the dielectric between the conductors responds with
an electric field $\vE$, and charge collects on the conductor surfaces. The
capacitance matrix relates the charges $\mathbf q=(q_1,\dots,q_n)$ to the applied
voltages $\mathbf v=(V_1,\dots,V_n)$ by $\mathbf q = \mat{C}\,\mathbf v$. Nothing
moves and nothing oscillates---this is a \emph{static} problem,
$\partial/\partial t=0$ (\cref{ch:maxwell}, \cref{sec:regimes}).
\end{physicsbox}

\begin{figure}[t]
\centering
\begin{tikzpicture}[scale=1.0]
  % dielectric background
  \fill[simBgBlue] (-3.6,-1.7) rectangle (3.6,1.7);
  \node[font=\scriptsize, text=simBlue] at (2.85,1.35) {dielectric $\varepsilon$};
  % conductor 1 (left plate, held at +V)
  \fill[simInk] (-2.7,0.55) rectangle (-0.4,0.95);
  \node[font=\scriptsize, text=white] at (-1.55,0.75) {conductor 1};
  \node[font=\scriptsize] at (-3.15,0.75) {$V_1$};
  % conductor 2 (right plate, held at 0)
  \fill[simInk] (-2.7,-0.95) rectangle (-0.4,-0.55);
  \node[font=\scriptsize, text=white] at (-1.55,-0.75) {conductor 2};
  \node[font=\scriptsize] at (-3.15,-0.75) {$V_2$};
  % field lines E = -grad V between the plates
  \foreach \x in {-2.4,-1.9,-1.4,-0.9}{
    \draw[flow, simRust, shorten <=1pt, shorten >=1pt] (\x,0.5)--(\x,-0.5);}
  \node[font=\scriptsize\itshape, text=simRust] at (-0.05,0.0) {$\vE=-\Grad V$};
  % a third, floating conductor on the right
  \fill[simInk] (1.3,-0.5) circle (0.45);
  \node[font=\scriptsize, text=white] at (1.3,-0.5) {3};
  \draw[flow, simRust, shorten >=1pt] (0.55,0.8) .. controls (1.0,0.4) .. (1.0,-0.2);
  \draw[flow, simRust, shorten >=1pt] (2.05,0.8) .. controls (1.6,0.4) .. (1.6,-0.2);
  % equipotential note
  \node[font=\scriptsize, text=simGray, align=center] at (2.5,-1.25)
    {each conductor:\\ an equipotential};
\end{tikzpicture}
\caption[The electrostatic setup]{The physical setup. A set of conductors
(dark) is embedded in a dielectric of permittivity $\varepsilon$ (blue). Each
conductor is a perfect equipotential held at a chosen voltage $V_i$; the
dielectric between them carries the electric field $\vE=-\Grad V$ (rust lines,
running from high potential to low and meeting every conductor at a right angle).
The simulation prescribes the voltages and computes how much charge collects on
each conductor per applied volt---the capacitance matrix. The field lines from the
upper plate that terminate on the floating conductor~3 are exactly the
\emph{mutual} coupling the off-diagonal entries of $\mat{C}$ measure.}
\label{fig:es-setup}
\end{figure}

The governing physics is the simplest in the book, and we derived it in full in
\cref{sec:electrostatic}. With $\vE=-\Grad V$ (a scalar potential, which makes the
static law $\Curl\vE=\mathbf0$ hold identically) and the dielectric constitutive
relation $\vD=\varepsilon\vE$, Gauss's law $\Div\vD=\rho$ collapses to the
generalized Poisson equation
\begin{equation}
  -\Div(\varepsilon\,\Grad V) = \rho .
  \label{eq:es-poisson-ch39}
\end{equation}
In the charge-free dielectric \emph{between} the conductors $\rho=0$, so the
equation Palace actually solves in the interior is the homogeneous Laplace
equation $-\Div(\varepsilon\Grad V)=0$; all the free charge lives on the conductor
surfaces, and it enters the problem through the boundary conditions, not the
right-hand side. The unknown $V$ is a scalar field, so it lives in the scalar
finite-element space $\Hone$ (continuous, nodal) of \cref{ch:functionspaces}---the
first slot of the de Rham complex, the natural home for a potential.

\subsection{The method of unit excitations}
\label{sec:unit-excitations}

The capacitance matrix has $n^2$ entries, but we do not solve $n^2$ problems to
find them. We solve $n$. The trick is the \emph{method of unit excitations}, and
it is the structural heart of the whole analysis.

The relation $\mathbf q=\mat{C}\,\mathbf v$ is linear. Its $j$th column,
$\mat{C}\,\mathbf e_j$, is the vector of charges produced when the voltage vector
is the $j$th standard basis vector $\mathbf e_j$---that is, when conductor $j$ is
held at \emph{one volt} and all the others are \emph{grounded} at zero. So to read
off column $j$ of $\mat{C}$ we set up exactly that excitation, solve the Laplace
problem for the resulting potential field $V_j$, and extract the charges. Repeat
for each conductor $j=1,\dots,n$ and we have all $n$ columns---the whole matrix
(\cref{fig:es-unit}). The $n$ solves produce a \emph{family} of potential fields
$\{V_1,\dots,V_n\}$, one per terminal, and that family is the raw material the
reduction stage turns into $\mat{C}$.

\begin{figure}[t]
\centering
\begin{tikzpicture}[scale=1.0, font=\footnotesize]
  % three little panels, one per unit excitation
  \foreach \k/\dx in {1/0, 2/4.2, 3/8.4}{
    \begin{scope}[shift={(\dx,0)}]
      \draw[simGray, semithick, rounded corners=2pt] (-0.55,-1.25) rectangle (3.0,1.35);
      \node[font=\scriptsize\bfseries] at (1.25,1.05) {excite conductor \k};
    \end{scope}}
  % panel 1: conductor 1 at 1V, others 0
  \begin{scope}[shift={(0,0)}]
    \fill[simRust] (0.0,0.45) rectangle (1.0,0.7);  \node[font=\scriptsize] at (-0.3,0.57){$1$};
    \fill[simInk] (0.0,-0.7) rectangle (1.0,-0.45); \node[font=\scriptsize] at (-0.3,-0.57){$0$};
    \fill[simInk] (1.7,-0.1) circle (0.22);          \node[font=\scriptsize] at (2.25,-0.1){$0$};
    \foreach \x in {0.2,0.5,0.8}{\draw[flow,simBlue,shorten <=1pt,shorten >=1pt](\x,0.4)--(\x,-0.4);}
    \node[font=\scriptsize\itshape] at (0.5,-1.02){$\to V_1$};
  \end{scope}
  % panel 2: conductor 2 at 1V
  \begin{scope}[shift={(4.2,0)}]
    \fill[simInk] (0.0,0.45) rectangle (1.0,0.7);   \node[font=\scriptsize] at (-0.3,0.57){$0$};
    \fill[simRust] (0.0,-0.7) rectangle (1.0,-0.45);\node[font=\scriptsize] at (-0.3,-0.57){$1$};
    \fill[simInk] (1.7,-0.1) circle (0.22);          \node[font=\scriptsize] at (2.25,-0.1){$0$};
    \foreach \x in {0.2,0.5,0.8}{\draw[flow,simBlue,shorten <=1pt,shorten >=1pt](\x,-0.4)--(\x,0.4);}
    \node[font=\scriptsize\itshape] at (0.5,-1.02){$\to V_2$};
  \end{scope}
  % panel 3: conductor 3 at 1V
  \begin{scope}[shift={(8.4,0)}]
    \fill[simInk] (0.0,0.45) rectangle (1.0,0.7);   \node[font=\scriptsize] at (-0.3,0.57){$0$};
    \fill[simInk] (0.0,-0.7) rectangle (1.0,-0.45); \node[font=\scriptsize] at (-0.3,-0.57){$0$};
    \fill[simRust] (1.5,-0.1) circle (0.24);         \node[font=\scriptsize] at (2.25,-0.1){$1$};
    \draw[flow,simBlue,shorten >=1pt](1.25,0.3) .. controls (0.9,0.0) .. (1.0,-0.05);
    \draw[flow,simBlue,shorten >=1pt](1.0,0.55) .. controls (0.7,0.2) .. (0.85,0.0);
    \node[font=\scriptsize\itshape] at (1.25,-1.02){$\to V_3$};
  \end{scope}
\end{tikzpicture}
\caption[The method of unit excitations]{The method of unit excitations on three
conductors. To recover column $j$ of the capacitance matrix we hold conductor $j$
at \emph{one volt} (rust) and ground all the others (dark, $0$), then solve the
Laplace problem $-\Div(\varepsilon\Grad V_j)=0$ for the resulting potential field
$V_j$ (blue field lines). Three excitations give three solved fields
$\{V_1,V_2,V_3\}$---the solution \emph{family}. Each solve has \emph{the same}
governing operator; only the prescribed boundary voltage moves. This is why the
analysis sits in the fixed-operator corner of \cref{ch:situating}.}
\label{fig:es-unit}
\end{figure}

\begin{keyidea}
The electrostatic analysis never assembles the capacitance matrix directly. It
exploits linearity: column $j$ of $\mat{C}$ is what you get by holding conductor
$j$ at unit voltage and the rest at ground, solving for the potential, and reading
off the charges. So $n$ conductors need $n$ solves, producing a \emph{family}
$\{V_1,\dots,V_n\}$ of potential fields. Each solve uses the \emph{same} discrete
operator and differs only in the prescribed boundary voltage---which is exactly the
shape that makes the next three sections possible.
\end{keyidea}

With the physical question and the method fixed, we can run the skeleton. The
electrostatic pipeline is
\[
  \texttt{config} \to \underbrace{\texttt{assemble}}_{\S\ref{sec:es-assemble}}
  \to \underbrace{\texttt{solve}}_{\S\ref{sec:es-solve}}
  \to \underbrace{\texttt{reduce}}_{\S\ref{sec:es-reduce}} \to \texttt{product},
\]
and the next three sections are those three stages, in order. We then post-process
the fields (\S\ref{sec:es-fields}) and step back to read the whole pipeline as the
anchor for Part~V (\S\ref{sec:es-anchor}).

\section{Assemble, once}
\label{sec:es-assemble}

The first stage turns the continuous physics of \cref{eq:es-poisson-ch39} into a
discrete operator. We did the mathematics in \cref{ch:weak} and \cref{ch:assembly};
here we only need its result and---crucially---\emph{when} it runs.

Multiplying the Laplace equation by a test function $v\in\Hone$ and integrating by
parts (the weak form, \cref{ch:weak}) turns the differential operator
$-\Div(\varepsilon\Grad\,\cdot\,)$ into the $\varepsilon$-weighted \emph{stiffness}
bilinear form
\begin{equation}
  a(u,v) \;=\; \int_\dom \varepsilon\,\Grad u \cdot \Grad v \dV ,
  \label{eq:es-bilinear-ch39}
\end{equation}
the single \code{weak\_form\_term} pairing the coefficient $\varepsilon$ with the
\emph{gradient} differential operator. Assembled over the mesh by the
assembly fold of \cref{ch:assembly}, this becomes a sparse, symmetric,
positive-definite matrix---the electrostatic stiffness operator $\Kop$. In the
framework's vocabulary the assembly is \code{fe\_assemble}, a fold that sums one
contribution per weak-form term; electrostatics is its simplest possible use,
because the term list has exactly one entry:
\begin{lstlisting}[style=l4style]
-- assemble the H1 stiffness operator from the single epsilon-weighted gradient term
space = h1_space cfg                          -- the H1 finite-element space (read-only)
terms = [ diffusion (permittivity cfg) ]      -- a one-term weak form: epsilon * grad
k     = fe_assemble space terms               -- fold the term list -> the stiffness K
\end{lstlisting}
The space and the term list are read-only construction-stratum data
(\cref{ch:strata}); $\Kop$ is a fresh value the fold returns. There is no mutation
and no loop we write---the loop is \emph{inside} \code{fe\_assemble}, over the term
list, and that list has length one.

\subsection{The boundary conditions enter by elimination}

The conductor voltages are essential (Dirichlet) boundary conditions
(\cref{ch:maxwell}, \S\ref{sec:bc}; \cref{ch:bc}): they constrain the value of $V$
on the conductor surfaces, so they must be built into the solution space rather
than added as a source. The mechanism is \emph{elimination}: the degrees of
freedom pinned by the boundary condition are removed from the unknown vector, and
their known values are moved to the right-hand side. The constrained operator
$\Kop$ is the stiffness matrix restricted to the free (interior) degrees of
freedom, and it is this reduced, still-symmetric-positive-definite matrix that the
solver receives. The free charge that Gauss's law puts on the conductor surfaces
never appears as a volume source; it is implicit in the eliminated boundary data.

\subsection{The load-bearing word is \emph{once}}

Here is the fact this section exists to make. The operator $\Kop$ is assembled
\emph{exactly once}, before any solving begins, and then held unchanged for the
entire analysis. Every one of the $n$ unit-excitation solves uses the same $\Kop$.
The driver makes this concrete: it builds the Krylov solver and hands it the
operator with a single \code{SetOperators} call, and that call sits \emph{outside}
the terminal loop.\footnote{In the Palace driver the operator is assembled by
\code{laplace\_op.GetStiffnessMatrix()} and captured by \code{ksp.SetOperators(*K,
*K)} \emph{before} the per-terminal \code{for} loop begins
(\code{palace/drivers/electrostaticsolver.cpp:30,36,59}). The placement of that one
call relative to that one loop is the entire structural signature of the
fixed-operator corner.}

\begin{pitfall}
The single most important line in the electrostatic driver is not a line of
arithmetic---it is the \emph{position} of the \code{SetOperators} call. It belongs
\emph{above} the terminal loop. An implementation that ``tidied'' it inside the
loop would reassemble (or at least re-factor and re-capture) the operator on every
terminal, doing $n$ times the setup work for no change in the answer. Worse, it
would erase the very property that distinguishes electrostatics from the driven
analysis of \cref{ch:driven}, where the operator genuinely \emph{does} change each
step and the call \emph{must} live inside the loop. Hoisting the operator is not an
optimization detail bolted onto the analysis; it \emph{is} the analysis's identity.
\end{pitfall}

\section{Solve: the fixed-operator map}
\label{sec:es-solve}

With $\Kop$ assembled and captured, the middle stage runs the family of solves.
This is the stage with all the structure, and it is the canonical occupant of the
first solve corner of \cref{ch:situating}: the \emph{fixed-operator map}.

\subsection{One operator, a fan of right-hand sides}

For each terminal $j$ the driver forms the right-hand side $\vb_j$ that encodes
``conductor $j$ at one volt, the rest grounded''---the eliminated boundary data of
the previous section, assembled into a vector. It then solves the linear system
\begin{equation}
  \Kop\,V_j \;=\; \vb_j
  \label{eq:es-system}
\end{equation}
for the potential field $V_j$. The operator $\Kop$ on the left is identical across
all $n$ systems; only the right-hand side $\vb_j$ changes. The $n$ systems are
therefore \emph{independent}---solving for $V_2$ needs nothing from $V_1$---so the
collection of solves is a \texttt{map} over the family of right-hand sides, not a
fold. They could be done in any order, or on $n$ separate machines at once
(\cref{fig:es-solvefamily}). This is the framework combinator \code{solve\_family},
and electrostatics is its defining witness:
\begin{lstlisting}[style=l4style]
-- the fixed-operator map: capture K ONCE, map the solver over the RHS family
solve_family :: LinOp -> [Tensor] -> [Tensor]
solve_family k bs = map (ksp_solve k) bs       -- k is captured once; reused per RHS

rhss = [ excitation cfg idx | idx <- terminal_sources cfg ]  -- one RHS per terminal
vs   = solve_family k rhss                       -- the solution family [V_1 .. V_n]
\end{lstlisting}
Read the body slowly. The operator \code{k} is a \emph{free variable} of the
mapped function \code{ksp\_solve k}; the \texttt{map} closes over it once and reuses
it for every element. That closure \emph{is} the hoist---the operator-capture made
structural by the type, so that no implementation can accidentally rebuild it
inside the loop. Each element solve \code{ksp\_solve k b\_j} is one preconditioned
conjugate-gradient solve of the SPD system \cref{eq:es-system}.

\begin{figure}[t]
\centering
\begin{tikzpicture}[font=\footnotesize, node distance=6mm]
  % the single fixed operator
  \node[opbox, minimum width=14mm, minimum height=12mm, fill=simBgBlue, draw=simBlue]
    (K) {$\Kop$};
  \node[font=\scriptsize\itshape, above=1.5mm of K] {assembled \emph{once}};
  % the fan of RHS and solves
  \foreach \j/\y in {1/2.2, 2/0.9, 3/-0.4, 4/-1.7}{
    \node[data, minimum width=9mm, right=20mm of K, yshift=(\y-0.25)*10mm]
      (b\j) {$\vb_{\j}$};
    \node[product, minimum width=10mm, right=14mm of b\j] (v\j) {$V_{\j}$};
    \draw[flow, shorten >=1pt] (K.east) .. controls +(0.9,0) .. (b\j.west);
    \draw[flow, shorten >=1pt] (b\j.east) -- node[above,font=\scriptsize]{\code{ksp\_solve}} (v\j.west);
  }
  \node[font=\scriptsize\itshape, text=simGreen, align=center]
    at ($(v4)+(0.1,-1.0)$) {independent solves:\\ embarrassingly parallel over terminals};
\end{tikzpicture}
\caption[\code{solve\_family} as a fixed-operator map]{The solve stage as a
\emph{fixed-operator map}. One operator $\Kop$ (blue, assembled once) feeds a fan
of independent solves, each with its own right-hand side $\vb_j$ (the $j$th unit
excitation) and producing its own potential field $V_j$. Because every solve
reuses the same $\Kop$ and none depends on another's result, the family is a
\texttt{map}: reorderable, parallelizable, ``embarrassingly parallel over
terminals.'' This is the combinator \code{solve\_family} of \cref{ch:combinators},
and the electrostatic terminal sweep is its defining witness. Contrast the driven
analysis (\cref{ch:driven}), where a \emph{fresh} operator feeds each solve.}
\label{fig:es-solvefamily}
\end{figure}

\subsection{The element solve is preconditioned CG}

Each \code{ksp\_solve} is a call to the conjugate-gradient method we built in
\cref{ch:cg}. CG is the right engine here for a specific reason: the constrained
stiffness $\Kop$ is symmetric positive-definite, and CG is the optimal Krylov
method for exactly that class (\cref{ch:cg}). In practice it is run with a
preconditioner---a cheap approximate inverse $\mat{P}^{-1}\approx\Kop^{-1}$ that
clusters the spectrum and cuts the iteration count, as developed in
\cref{ch:precond}. The convergence is the geometric decay of the CG error in the
$\Kop$-norm; \cref{fig:es-cg} shows a representative residual history for one
terminal solve, with and without preconditioning. The point for \emph{this} chapter
is only that the element solve is a self-contained, well-understood SPD solve---a
single piece of Part~IV machinery---and that the family runs $n$ of them over the
one operator.

\begin{figure}[t]
\centering
\begin{tikzpicture}
\begin{axis}[
    width=0.74\linewidth, height=52mm,
    xlabel={CG iteration $k$}, ylabel={relative residual $\norm{\vr_k}/\norm{\vb}$},
    xlabel style={font=\scriptsize}, ylabel style={font=\scriptsize},
    xmin=0, xmax=60, ymin=1e-8, ymax=2,
    ymode=log, ytick={1,1e-2,1e-4,1e-6,1e-8},
    xtick={0,20,40,60},
    yticklabel style={font=\scriptsize}, xticklabel style={font=\scriptsize},
    legend style={font=\scriptsize, draw=simGray!50, at={(0.98,0.98)},
      anchor=north east},
    every axis plot/.append style={very thick},
    axis lines=left,
  ]
  % unpreconditioned: slow geometric decay
  \addplot[simRust, domain=0:60, samples=61] {exp(-0.26*x)};
  \addlegendentry{plain CG}
  % preconditioned: fast geometric decay
  \addplot[simBlue, domain=0:18, samples=19] {exp(-0.95*x)};
  \addlegendentry{preconditioned CG}
  % convergence tolerance line
  \addplot[simGray, densely dashed, domain=0:60, samples=2] {1e-6};
  \addlegendentry{tolerance}
\end{axis}
\end{tikzpicture}
\caption[CG convergence for one terminal solve]{A representative convergence
history for a single unit-excitation solve $\Kop V_j=\vb_j$. The relative residual
falls geometrically (a straight line on the log axis) until it crosses the
convergence tolerance (dashed). A good preconditioner (\cref{ch:precond}, blue)
clusters the spectrum of $\Kop$ and reaches tolerance in a fraction of the
iterations of plain CG (\cref{ch:cg}, rust). Because the operator is fixed across
the whole family, the same preconditioner is built once and reused for every
terminal solve---another dividend of the hoist.}
\label{fig:es-cg}
\end{figure}

\begin{notationbox}
The fixed operator pays a second dividend beyond avoiding reassembly. Because
$\Kop$ does not change, any \emph{preconditioner} built from it---an incomplete
factorization, a multigrid hierarchy (\cref{ch:precond})---is also built once and
reused across all $n$ solves. The expensive setup is amortized over the whole
family. A direct factorization would amortize even more dramatically: factor
$\Kop=\mat{L}\mat{L}^{\!\top}$ once, then each terminal is a pair of cheap
triangular solves. ``Assemble once'' quietly licenses ``factor once'' and
``precondition once''---the entire economic case for the fixed-operator corner.
\end{notationbox}

\section{Reduce: the symmetric Gram}
\label{sec:es-reduce}

The solve stage has handed us the family $\{V_1,\dots,V_n\}$. The reduction stage
turns it into the capacitance matrix. We met this reduction in full generality in
\cref{ch:reductions}; here it appears in its simplest and most important instance.

\subsection{Capacitance is an operator-weighted Gram}

The $(i,j)$ entry of the capacitance matrix is the operator-weighted bilinear form
of two family members,
\begin{equation}
  \boxed{\;C_{ij} \;=\; V_j^{\!\top}\,\Kop\,V_i\;}
  \label{eq:cap-gram-ch39}
\end{equation}
where $\Kop$ is the $\varepsilon$-weighted energy operator (the discrete form of
the electric energy density $\tfrac12\varepsilon\abs{\vE}^2$ of \cref{ch:maxwell},
eq.~\eqref{eq:energy-density}). This is the \emph{symmetric Gram reduction}
$G_{ij}=w(i,j)\,V_j^{\!\top}\Kop V_i$ of \cref{ch:reductions}, with the weight set
to unity, $w(i,j)=1$. The framework verb is \code{gram\_reduce}, and capacitance is
its $w\equiv1$ specialization:
\begin{lstlisting}[style=l4style]
-- the symmetric Gram reduction; capacitance is the unit-weight (w = 1) member
gram_reduce :: LinOp -> [Tensor] -> (Int -> Int -> Scalar) -> Matrix
gram_reduce k vs w =
  symmetric_from_upper                              -- mirror lower triangle from upper
    [ [ w i j * entry k vs i j | j <- [i .. m-1] ]  -- map over upper-triangle pairs
      | i <- [0 .. m-1] ]
  where
    m = length vs
    entry k vs i j
      | i == j    = matrix_weighted_norm (vs!!i) k  -- diagonal: V_i^T K V_i
      | otherwise = bilinear_form (vs!!j) k (vs!!i) -- off-diag:  V_j^T K V_i

capacitance = gram_reduce k vs (\i j -> 1)          -- w = 1: unit voltage excitation
\end{lstlisting}
Two structural facts are built into this body, and both are load-bearing for
capacitance specifically. First, the grid of entries is a \texttt{map} over family
pairs with no state threaded between them---the embarrassingly parallel pattern of
\cref{ch:reductions}: no entry needs another's value. Second, the matrix is
\emph{symmetric by construction}. Because $\Kop$ is symmetric,
$V_j^{\!\top}\Kop V_i = V_i^{\!\top}\Kop V_j$, so the verb computes only the upper
triangle and mirrors the lower (\code{symmetric\_from\_upper}). For an $n$-conductor
problem this is $\binom{n}{2}+n$ bilinear evaluations rather than $n^2$---the
symmetry is exploited, not merely observed (\cref{fig:es-gram}).

\begin{figure}[t]
\centering
\begin{tikzpicture}[font=\footnotesize]
  % the symmetric Gram grid
  \foreach \i in {0,1,2}{
    \foreach \j in {0,1,2}{
      \node[draw=simGray, semithick,
        fill={\ifnum\i=\j simBgRust\else simBgBlue\fi},
        minimum size=8mm, inner sep=0pt]
        (g\i\j) at (\j*0.84, -\i*0.84) {};}}
  % labels on entries
  \node[font=\scriptsize] at (0,0)        {$C_{11}$};
  \node[font=\scriptsize] at (0.84,-0.84) {$C_{22}$};
  \node[font=\scriptsize] at (1.68,-1.68) {$C_{33}$};
  \node[font=\scriptsize, simBlue] at (0.84,0)    {$C_{12}$};
  \node[font=\scriptsize, simBlue] at (0,-0.84)   {$C_{21}$};
  % mirror line
  \draw[densely dashed, simRust, semithick] (-0.55,0.55) -- (2.23,-2.23);
  \node[font=\scriptsize, text=simRust] at (2.0,0.4) {mirror};
  % computed-half brace
  \node[font=\scriptsize, align=center, text=simBlue] at (3.6,-0.1)
    {upper triangle\\ \emph{computed}:\\ $V_j^{\!\top}\Kop V_i$};
  \node[font=\scriptsize, align=center, text=simRust] at (3.6,-1.7)
    {lower triangle\\ \emph{mirrored}:\\ $C_{ji}=C_{ij}$};
  % diagonal = energy callout
  \node[font=\scriptsize, align=center, text=simRust] at (1.0,-2.7)
    {diagonal $=$ twice the stored field energy:\ $C_{ii}=V_i^{\!\top}\Kop V_i = 2U_e$};
\end{tikzpicture}
\caption[The capacitance Gram]{The capacitance matrix as a symmetric Gram. Each
entry is the operator-weighted bilinear form $C_{ij}=V_j^{\!\top}\Kop V_i$ of two
solved potential fields. The diagonal (rust) is a \emph{self}-pairing,
$C_{ii}=V_i^{\!\top}\Kop V_i$, which equals \emph{twice} the field energy stored
when conductor $i$ alone is at unit voltage---the self-capacitance. The
off-diagonal (blue) is a \emph{cross}-pairing of two different fields, the mutual
capacitance. Because $\Kop$ is symmetric, $C_{ij}=C_{ji}$: only the upper triangle
is computed and the lower is mirrored across the dashed diagonal. The symmetry is
structural, a consequence of the symmetric energy operator, not a coincidence to be
checked afterward.}
\label{fig:es-gram}
\end{figure}

\subsection{Self versus mutual; the energy diagonal}

The two halves of the matrix carry two distinct physical meanings.

The \emph{diagonal} $C_{ii}=V_i^{\!\top}\Kop V_i$ is a self-pairing. Applying the
energy operator $\Kop$ to the field $V_i$ and pairing the result back with $V_i$
yields exactly twice the electric field energy stored when conductor $i$ alone is
held at unit voltage. This is the \emph{energy form} of capacitance,
\begin{equation}
  C_{ii} \;=\; \frac{2U_e}{V_i^2}\Bigg|_{V_i=1} \;=\; 2U_e ,
  \qquad U_e=\tfrac12\!\int_\dom \varepsilon\abs{\vE}^2\dV ,
  \label{eq:cap-energy}
\end{equation}
the relation \cref{ch:maxwell} promised (eq.~\eqref{eq:energy-density}, $C=2U/V^2$),
specialized to the unit voltage we excite with. It is the same number a
field-energy integration would produce---which is exactly the cross-check of
Worked example~\ref{wex:es-energy} below. (This energy formulation is the one
Palace's driver actually computes, following the COMSOL AC/DC manual's
$C_{ii}=2U_e/V_i^2$ recipe.)

The \emph{off-diagonal} $C_{ij}=V_j^{\!\top}\Kop V_i$ is a cross-pairing of two
different fields, the \emph{mutual} capacitance: the charge induced on conductor
$i$ by a unit voltage on conductor $j$. A robust and slightly surprising fact is
that the off-diagonal of a Maxwell capacitance matrix is always \emph{negative}.
Raising conductor $j$'s voltage draws charge of the \emph{opposite} sign onto a
neighbouring grounded conductor through their shared field, so in
$q_i=C_{ii}V_i+\sum_{j\neq i}C_{ij}V_j$ the coupling term $C_{ij}V_j$ must reduce
$q_i$ as $V_j$ rises---hence $C_{ij}<0$. The matrix is symmetric, $C_{ij}=C_{ji}$,
a reciprocity the Gram construction makes automatic: the charge conductor $j$
induces on conductor $i$ per volt equals the charge $i$ induces on $j$ per volt,
because the energy operator is symmetric.

\begin{figure}[t]
\centering
\begin{tikzpicture}[font=\footnotesize, scale=1.0]
  % a 3x3 matrix with self (diag) vs mutual (off-diag) highlighted
  \matrix[matrix of math nodes, left delimiter={(}, right delimiter={)},
    nodes={minimum size=8mm, anchor=center},
    column sep=2pt, row sep=2pt] (C) at (0,0)
  {
    {\color{simRust}C_{11}} & {\color{simBlue}C_{12}} & {\color{simBlue}C_{13}} \\
    {\color{simBlue}C_{21}} & {\color{simRust}C_{22}} & {\color{simBlue}C_{23}} \\
    {\color{simBlue}C_{31}} & {\color{simBlue}C_{32}} & {\color{simRust}C_{33}} \\
  };
  % annotations
  \node[font=\scriptsize, text=simRust, align=left, anchor=west] at (2.3,0.85)
    {\textbf{diagonal} $C_{ii}>0$:\\ \emph{self}-capacitance\\ ($=2U_e$, the stored energy)};
  \node[font=\scriptsize, text=simBlue, align=left, anchor=west] at (2.3,-0.7)
    {\textbf{off-diagonal} $C_{ij}<0$:\\ \emph{mutual} capacitance\\ ($C_{ij}=C_{ji}$, reciprocity)};
\end{tikzpicture}
\caption[Self versus mutual capacitance]{The structure of the Maxwell capacitance
matrix. The \emph{diagonal} entries (rust) are the self-capacitances; each is
positive and equals twice the field energy stored with that conductor alone at unit
voltage. The \emph{off-diagonal} entries (blue) are the mutual capacitances; each is
\emph{negative}---a unit voltage on conductor $j$ induces opposite-sign charge on
conductor $i$---and the matrix is symmetric, $C_{ij}=C_{ji}$, by the reciprocity the
symmetric Gram construction guarantees. Inverting $\mat{C}$ gives the alternate
\emph{Maxwell} form relating the conductor potentials to their charges,
$\mathbf v=\mat{C}^{-1}\mathbf q$.}
\label{fig:es-selfmutual}
\end{figure}

\begin{keyidea}
The capacitance matrix is the symmetric Gram reduction
$C_{ij}=V_j^{\!\top}\Kop V_i$ of the per-terminal potential family against the
energy operator $\Kop$, with unit weight. Its diagonal is twice the stored field
energy (self-capacitance, always positive); its off-diagonal is the mutual
capacitance (always negative); and it is symmetric by construction because $\Kop$
is. Capacitance and inductance (\cref{ch:magnetostatics}) call the \emph{same}
verb, \code{gram\_reduce}, and differ only in the weight $w$---one verb, two
products.
\end{keyidea}

\subsection{The product, and its inverse}

The verb returns the $n\times n$ matrix $\mat{C}$. This is the physical product the
user ran the analysis to compute. Palace additionally forms the inverse
$\mat{C}^{-1}$ by a dense LAPACK factorization---a small, cheap operation on the
$n\times n$ result, where $n$ is the handful of conductors, not the millions of
mesh degrees of freedom. The inverse is the \emph{other} Maxwell form: where
$\mathbf q=\mat{C}\,\mathbf v$ gives charges from voltages, $\mathbf
v=\mat{C}^{-1}\mathbf q$ gives the voltages that a prescribed set of charges would
induce. The inversion is a downstream consumer of the reduction's output, not part
of the Gram reduction itself.

\section{Field post-processing}
\label{sec:es-fields}

Before assembling the whole pipeline, we account for one more thing the analysis
produces: the fields themselves. The capacitance matrix is the headline number, but
each solved potential $V_j$ also yields a visualizable electric field and a stored
energy, and these are computed as cheap post-processes of the family.

The electric field of excitation $j$ is the (negative) gradient of its potential,
\begin{equation}
  \vE_j \;=\; -\Grad V_j ,
  \label{eq:es-Efield}
\end{equation}
computed by applying the discrete gradient operator to the solved nodal potential
(the same $\Grad$ matrix the assembly produced). The energy stored in that field is
\begin{equation}
  U_e^{(j)} \;=\; \tfrac12\!\int_\dom \varepsilon\,\abs{\vE_j}^2 \dV
            \;=\; \tfrac12\,V_j^{\!\top}\Kop V_j ,
  \label{eq:es-energy}
\end{equation}
and the second equality is the bridge that makes the diagonal of $\mat{C}$ an energy:
the field-energy integral is \emph{exactly} half the Gram diagonal entry $C_{jj}$.
This is no coincidence---the energy operator $\Kop$ is, by construction
(\cref{ch:maxwell}, eq.~\eqref{eq:energy-density}), the discrete quadratic form
whose value on a field is twice its stored energy. The same observation lets the
\emph{domain-energy} reduction of \cref{ch:reductions} report \emph{where} the energy
sits---what fraction lives in each dielectric region (the participation factors of
\cref{ch:maxwell}, \S\ref{sec:Q})---by restricting the energy operator to each
domain. The energy is also the thread back to Poynting's theorem (\cref{ch:maxwell},
\S\ref{sec:energy}): in the static case the Poynting flux is zero and all the energy
is the stored electric energy \cref{eq:es-energy}, the term that the capacitance
reports.

\section{The anchor: the whole pipeline, and what it anchors}
\label{sec:es-anchor}

We have now traced every stage. It is time to stand back and see the analysis
whole, because its \emph{shape}---not its electrostatics---is what the rest of
Part~V is built on.

\subsection{The pipeline in one line}

The three stages compose into a single expression. At the framework level the whole
electrostatic feature is
\begin{lstlisting}[style=l4style]
electrostatic = gram_reduce . solve_family . fe_assemble
\end{lstlisting}
read right to left: \code{fe\_assemble} builds the operator once, \code{solve\_family}
maps the family of solves over it, and \code{gram\_reduce} collapses the solved
family to the capacitance matrix. Spelled out with the data flowing through it,
\begin{lstlisting}[style=l4style]
-- inputs = config; output = the capacitance matrix (the physical product)
electrostatic :: ElectrostaticConfig -> CapacitanceMatrix
electrostatic cfg =
  let space = h1_space cfg                          -- the H1 finite-element space
      k     = fe_assemble space [ diffusion (permittivity cfg) ]  -- (1) assemble K ONCE
      rhss  = [ excitation cfg idx | idx <- terminal_sources cfg ] -- per-terminal RHS family
      vs    = solve_family k rhss                    -- (2) fixed-operator per-terminal map
  in  gram_reduce k vs (\i j -> 1)                   -- (3) C_ij = V_j^T K V_i, w = 1
\end{lstlisting}
\Cref{fig:es-pipeline} draws it as a labeled diagram. This figure is the one to
remember: every later modality chapter is, quite literally, a redrawing of it with
one or two boxes replaced.

\begin{figure}[p]
\centering
\resizebox{\linewidth}{!}{%
\begin{tikzpicture}[node distance=7mm and 13mm, font=\footnotesize]
  % the pipeline, left to right
  \node[data, align=center, text width=20mm] (cfg)
    {config\\[1pt]\footnotesize $\varepsilon$, mesh,\\terminals};
  \node[stage, right=of cfg, align=center, text width=22mm] (asm)
    {\textbf{assemble}\\[1pt]\footnotesize \code{fe\_assemble}\\(once)};
  \node[stage, right=of asm, align=center, text width=24mm, fill=simBgGreen, draw=simGreen]
    (solve) {\textbf{solve}\\[1pt]\footnotesize \code{solve\_family}\\(per terminal)};
  \node[stage, right=of solve, align=center, text width=22mm, fill=simBgGold, draw=simGold]
    (red) {\textbf{reduce}\\[1pt]\footnotesize \code{gram\_reduce}\\($w{=}1$)};
  \node[product, right=of red, align=center, text width=18mm] (prod)
    {$\mat{C}$\\[1pt]\footnotesize capacitance\\matrix};
  % flow edges with the carried data labelled
  \draw[flow] (cfg) -- (asm);
  \draw[flow] (asm) -- node[above,font=\scriptsize]{$\Kop$} (solve);
  \draw[flow] (solve) -- node[above,font=\scriptsize]{$\{V_j\}$} (red);
  \draw[flow] (red) -- (prod);
  % the shared-operator capture: K also feeds the reduction
  \draw[depedge, shorten >=1pt] (asm.south) ++(0,-1mm)
    .. controls +(0,-9mm) and +(0,-9mm) .. (red.south)
    node[midway, below, font=\scriptsize, text=simBlue]
      {the \emph{same} $\Kop$ weights the Gram};
  % corner / shape annotations under each stage
  \node[font=\scriptsize\itshape, text=simGreen, below=10mm of solve, align=center]
    {fixed-operator map\\(\cref{ch:situating}, corner 1)};
  \node[font=\scriptsize\itshape, text=simGold, above=2mm of red, align=center]
    {rank-2 symmetric Gram};
\end{tikzpicture}%
}
\caption[The electrostatic pipeline---the anchor diagram]{The complete
electrostatic pipeline, $\texttt{electrostatic} = \code{gram\_reduce}\circ
\code{solve\_family}\circ\code{fe\_assemble}$. A configuration of permittivity,
mesh, and terminals enters; \code{fe\_assemble} builds the stiffness operator
$\Kop$ \emph{once}; \code{solve\_family} maps the per-terminal solves over that
fixed operator (the fixed-operator corner of \cref{ch:situating}), producing the
potential family $\{V_j\}$; and \code{gram\_reduce} pairs that family against
$\Kop$ into the capacitance matrix $\mat{C}$ (the rank-2 symmetric Gram). The blue
back-edge marks a feature special to this analysis: the \emph{same} operator
$\Kop$ that drove the solves also weights the reduction. \textbf{This diagram is
the anchor.} Every other analysis in Part~V is this pipeline with the FE space, the
solve corner, and/or the reduction swapped---and the swaps are small and
well-localized.}
\label{fig:es-pipeline}
\end{figure}

\subsection{Every other analysis is this, with one stage swapped}

Here is the payoff that earns this chapter its length. The electrostatic pipeline
of \cref{fig:es-pipeline} is the \emph{template} for all six analyses. Each of the
others is obtained by swapping one or two boxes---and nothing else:

\begin{itemize}
  \item \textbf{Magnetostatics} (\cref{ch:magnetostatics}) swaps the
  \emph{function space} ($\Hone\to\Hcurl$, scalar potential $\to$ vector potential)
  and the \emph{Gram weight} ($w=1\to w=1/(I_iI_j)$, the current normalization). The
  solve corner is unchanged---it is still a fixed-operator map---and the reduction
  is still \code{gram\_reduce}. Two cells move; the skeleton is identical. It is the
  closest sibling, almost line for line the same analysis.
  \item \textbf{The driven analysis} (\cref{ch:driven}) keeps the map but swaps the
  \emph{solve corner} from fixed to operator-varying: the system operator
  $\mat{A}(\omega)=\Kop+i\omega\Cop-\omega^2\Mop$ depends on frequency, so it is
  rebuilt \emph{inside} the sweep---the \code{SetOperators} call moves \emph{into}
  the loop, the exact negation of \S\ref{sec:es-assemble}'s hoist. The reduction
  also swaps, from the symmetric Gram to the port-projection
  (\cref{ch:reductions}).
  \item \textbf{The eigenmode analysis} (\cref{ch:eigenmodes}) swaps the
  \emph{solve corner} for the opaque black box: there is no loop we write at all---the
  assembled pencil $(\Kop,\Mop)$ goes to a library eigensolver with one call. The
  reduction swaps to the rank-1 $(f,Q)$ table.
  \item \textbf{The transient analysis} (\cref{ch:transient}) swaps the
  \emph{map for a fold}: each time step begins from the state the previous step
  produced, so the independent family becomes a chained, sequential evolution.
  \item \textbf{The boundary-mode analysis} (\cref{ch:waveports}) is the eigenmode
  swap, performed on a two-dimensional cross-section instead of the full volume.
\end{itemize}

\begin{table}[t]
\centering
\caption[The modalities as swaps of the anchor]{Every Part~V modality as a swap of
the electrostatic anchor (\cref{fig:es-pipeline}). The columns are the stages of the
skeleton; a dash means ``unchanged from electrostatics.'' Reading down a column shows
which stage each analysis perturbs; reading across a row shows how few cells move.
Magnetostatics is the closest sibling (space and weight only); the others each swap
the solve corner, and most also the reduction. This table is the map of Part~V.}
\label{tab:es-swaps}
\small
\setlength{\tabcolsep}{4.5pt}
\renewcommand{\arraystretch}{1.3}
\begin{tabular}{@{}lllll@{}}
\toprule
Analysis & FE space & Assemble & Solve corner & Reduction \\
\midrule
\textbf{Electrostatic} (anchor) & $\Hone$ & $\Kop$ once & fixed-operator map & Gram, $w{=}1$ \\
Magnetostatic & $\Hcurl$ & --- & --- & Gram, $w{=}\tfrac{1}{I_iI_j}$ \\
Driven & $\Hcurl$ & $\mat{A}(\omega)$ in loop & operator-varying map & projection \\
Eigenmode & $\Hcurl$ & pencil once & opaque black box & table $(f,Q)$ \\
Transient & $\Hcurl$ & --- & state-threaded fold & table (history) \\
Boundary mode & $\Hcurl{\oplus}\Hone$ & 2-D pencil & opaque black box & table (modes) \\
\bottomrule
\end{tabular}
\end{table}

\Cref{tab:es-swaps} lays the swaps out as a table. The dashes are the point: most
cells of most rows are ``unchanged from electrostatics.'' This is why the book can
afford to teach electrostatics once, here, in complete detail, and then present
each remaining modality as a short, well-localized set of differences from it. The
full end-to-end trace---this chapter's worked example below, but on a real
three-dimensional device---is revisited as the capstone of Part~V in
\cref{ch:fulltrace}, where we run the entire pipeline against actual machinery and
check the answer.

\begin{keyidea}
The electrostatic pipeline $\code{gram\_reduce}\circ\code{solve\_family}\circ
\code{fe\_assemble}$ is the \textbf{anchor} for all of Part~V: assemble the operator
\emph{once}, \texttt{map} the solve over the excitation family (the fixed-operator
corner), and Gram-reduce the solved family to the product. Every other analysis is
this template with the FE space, the solve corner, and/or the reduction swapped---
magnetostatics swaps the space and the Gram weight, the driven analysis swaps to an
operator-varying map and a projection, eigenmode swaps the solve to a black box,
transient swaps the map for a fold. Learn this pipeline and you have learned the
spine of every analysis Palace runs.
\end{keyidea}

\section{Two worked examples}
\label{sec:es-worked}

The pipeline is general; we now run it end to end on a problem small enough to do by
hand. The first example traces the whole analysis---assemble, solve the two unit
excitations, Gram-reduce, mirror, interpret---and recovers the textbook
parallel-plate formula. The second checks that the Gram diagonal really is the field
energy, closing the loop on \cref{eq:es-energy}.

\begin{workedexample}{A two-conductor capacitance, end to end}{es-twocond}
Two conductors form a parallel-plate-like capacitor: an upper plate (conductor~1)
and a lower plate (conductor~2), with a uniform dielectric between them. We discretize
the gap so coarsely that the free potential reduces to a \emph{single} interior
degree of freedom $u$---the potential at the midplane---which is all we need to see
every stage of the pipeline turn.

\medskip
\emph{Assemble (once).} On this one-degree-of-freedom mesh the $\varepsilon$-weighted
stiffness form \cref{eq:es-bilinear-ch39} assembles, before elimination, into a
$3\times3$ matrix over the nodes (plate~1, midplane, plate~2). Writing
$g=\varepsilon A/(L/2)$ for the conductance of each half-gap (plate area $A$,
half-gap $L/2$), the unconstrained stiffness is
\[
  \Kop_{\text{full}} =
  \begin{pmatrix} g & -g & 0 \\ -g & 2g & -g \\ 0 & -g & g \end{pmatrix}.
\]
The plate nodes are Dirichlet (essential) degrees of freedom; eliminating them
leaves the single free interior unknown $u$ with the scalar constrained operator
$\Kop=2g$ and a right-hand side built from the prescribed plate voltages. This is
the assemble stage, run \emph{once}; the same $\Kop=2g$ serves both solves below.

\medskip
\emph{Solve the family (fixed-operator map).} Two unit excitations, two solves of
$\Kop u=b$ with the \emph{same} $\Kop=2g$:
\begin{itemize}[nosep]
  \item \textbf{Excitation 1:} plate~1 at $1$, plate~2 at $0$. Elimination gives the
  right-hand side $b_1=g\cdot1+g\cdot0=g$, so $u=b_1/\Kop=g/(2g)=\tfrac12$. The solved
  field $V_1$ has plate~1 at $1$, the midplane at $\tfrac12$, plate~2 at $0$---a linear
  ramp, the discrete image of the parallel-plate solution of \cref{wex:plates1d}.
  \item \textbf{Excitation 2:} plate~2 at $1$, plate~1 at $0$. By symmetry $u=\tfrac12$
  again, and $V_2$ ramps the other way: plate~2 at $1$, midplane at $\tfrac12$,
  plate~1 at $0$.
\end{itemize}
The two solves reused one operator; neither needed the other. That is the
fixed-operator map in miniature.

\medskip
\emph{Reduce (symmetric Gram).} Now form $C_{ij}=V_j^{\!\top}\Kop_{\text{full}}V_i$
on the full nodal fields (the energy operator pairs over all nodes). With
$V_1=(1,\tfrac12,0)$ and $V_2=(0,\tfrac12,1)$ we compute the upper triangle:
\[
  C_{11}=V_1^{\!\top}\Kop_{\text{full}}V_1
       = g\!\left[(1{-}\tfrac12)^2+(\tfrac12{-}0)^2\right] = g\cdot\tfrac12 = \tfrac{g}{2},
\]
and by the same geometry $C_{22}=\tfrac{g}{2}$. The off-diagonal is
\[
  C_{12}=V_2^{\!\top}\Kop_{\text{full}}V_1
       = g\bigl[(0{-}\tfrac12)(1{-}\tfrac12)+(\tfrac12{-}1)(\tfrac12{-}0)\bigr]
       = g\bigl[-\tfrac14-\tfrac14\bigr] = -\tfrac{g}{2}.
\]
We do \emph{not} recompute $C_{21}$: the symmetric Gram mirrors it,
$C_{21}=C_{12}=-\tfrac{g}{2}$. Assembling,
\[
  \mat{C} =
  \begin{pmatrix} \tfrac{g}{2} & -\tfrac{g}{2} \\[2pt] -\tfrac{g}{2} & \tfrac{g}{2} \end{pmatrix}
  = \frac{g}{2}\begin{pmatrix} 1 & -1 \\ -1 & 1 \end{pmatrix},
  \qquad g=\frac{\varepsilon A}{L/2}=\frac{2\varepsilon A}{L}.
\]

\medskip
\emph{Interpret.} The off-diagonal is negative, as every Maxwell capacitance matrix
demands: raising plate~2 induces opposite-sign charge on plate~1. The matrix is
symmetric, $C_{12}=C_{21}$, by reciprocity. And the physically meaningful
two-terminal capacitance---the charge that flows between the plates per volt of
\emph{difference}---is read from the matrix as $C_{\text{cap}}=-C_{12}=g/2=\varepsilon
A/L$, recovering the textbook parallel-plate result $C=\varepsilon A/L$ of
\cref{wex:plates1d}. The whole pipeline---assemble once, solve the two excitations,
Gram-reduce, mirror---has run in a page of arithmetic and produced the right answer.
On a fine three-dimensional mesh the only thing that changes is that the solves are
millions of degrees of freedom handled by CG (\cref{ch:cg}) instead of one division;
the \emph{structure} is exactly what we just did.
\end{workedexample}

\begin{workedexample}{The energy form agrees with the Gram diagonal}{es-energy}
We check the claim of \cref{eq:es-energy}: the diagonal Gram entry $C_{11}$ equals
twice the stored field energy of excitation~1, computed independently from the field.

\medskip
\emph{The Gram diagonal.} From \cref{wex:es-twocond}, $C_{11}=g/2$ with $g=2\varepsilon
A/L$, so $C_{11}=\varepsilon A/L$.

\emph{The field energy, computed directly.} Excitation~1 produced the field $V_1$
ramping linearly from $1$ at plate~1 to $0$ at plate~2 across the gap of width $L$,
so its electric field is uniform, $\vE_1=-\Grad V_1$ with magnitude
$\abs{\vE_1}=1/L$ (a unit potential drop over a distance $L$). The stored electric
energy is the energy density $\tfrac12\varepsilon\abs{\vE_1}^2$ integrated over the
gap volume $A\cdot L$:
\[
  U_e^{(1)} = \tfrac12\!\int_\dom \varepsilon\abs{\vE_1}^2\dV
            = \tfrac12\,\varepsilon\,\frac{1}{L^2}\,(A\,L)
            = \frac{\varepsilon A}{2L}.
\]

\emph{Compare.} Twice the energy is $2U_e^{(1)}=\varepsilon A/L$, which is
\emph{exactly} $C_{11}$. The Gram diagonal and the field-energy integral agree, as
\cref{eq:es-energy} promised:
\[
  C_{11} \;=\; V_1^{\!\top}\Kop V_1 \;=\; 2U_e^{(1)} \;=\; \frac{\varepsilon A}{L}.
\]
This is not a numerical coincidence. The energy operator $\Kop$ \emph{is}, by
construction, the discrete quadratic form whose value on a field is twice its stored
energy (\cref{ch:maxwell}, eq.~\eqref{eq:energy-density}); so the self-pairing
$V_i^{\!\top}\Kop V_i$ and the energy integral $2U_e$ are two evaluations of the same
quadratic form. The agreement is the reason the diagonal of the capacitance matrix
\emph{is} a capacitance: it measures the field energy a conductor stores per unit
voltage, which is exactly what capacitance means.
\end{workedexample}

\summarysection
The electrostatic analysis computes the capacitance matrix $\mat{C}$ of a system of
$n$ conductors in a dielectric. It exploits linearity through the \emph{method of
unit excitations}: hold conductor $j$ at one volt, ground the rest, solve the Laplace
problem $-\Div(\varepsilon\Grad V_j)=0$ for the potential, and repeat for each
conductor, producing a \emph{family} $\{V_1,\dots,V_n\}$ of potential fields. The
pipeline runs the skeleton of \cref{ch:situating} in its simplest form. \emph{Assemble}
($\S\ref{sec:es-assemble}$): the $\varepsilon$-weighted stiffness operator $\Kop$ on
the scalar space $\Hone$, with the terminal voltages eliminated as essential boundary
conditions, built \emph{once} and held---the load-bearing hoist. \emph{Solve}
($\S\ref{sec:es-solve}$): the fixed-operator map \code{solve\_family}, $n$ independent
preconditioned-CG solves $\Kop V_j=\vb_j$ reusing the one operator---the defining
witness of the first solve corner. \emph{Reduce} ($\S\ref{sec:es-reduce}$): the
symmetric Gram \code{gram\_reduce} with unit weight, $C_{ij}=V_j^{\!\top}\Kop V_i$,
whose diagonal is twice the stored field energy (self-capacitance, positive) and whose
off-diagonal is the mutual capacitance (negative), symmetric by construction. The whole
feature is the one-line composition
$\texttt{electrostatic}=\code{gram\_reduce}\circ\code{solve\_family}\circ
\code{fe\_assemble}$. This pipeline is the \textbf{anchor} of Part~V: every other
modality is this diagram (\cref{fig:es-pipeline}) with the FE space, the solve corner,
and/or the reduction swapped (\cref{tab:es-swaps}).

\furtherreading
For the electrostatics itself---capacitance, the energy method $C=2U/V^2$, and the
sign and reciprocity of the Maxwell capacitance matrix---Griffiths~\citep{griffiths2017}
is the standard undergraduate source, and Jin's finite-element electromagnetics
text~\citep{jin2014} treats the discretization, the $\Hone$ stiffness assembly, and the
energy-form capacitance extraction in the exact form Palace uses. Pozar~\citep{pozar2011}
places the capacitance matrix in its microwave-engineering context, alongside the
inductance matrix it pairs with (\cref{ch:magnetostatics}) and the S-parameters of the
driven analysis (\cref{ch:driven}). The fixed-operator solve corner and the symmetric
Gram reduction are this book's own framing, developed in \cref{ch:situating} and
\cref{ch:reductions}; the numerical engine behind each element solve is the conjugate
gradient method of \cref{ch:cg} with the preconditioners of \cref{ch:precond}.

\exercisessection
\begin{enumerate}[nosep]
  \item \textbf{Why $n$ solves, not $n^2$.} The capacitance matrix has $n^2$ entries.
  Explain, using the linearity of $\mathbf q=\mat{C}\,\mathbf v$, why the analysis
  needs only $n$ solves to fill them all. What is each solve's right-hand side, and
  which column of $\mat{C}$ does it produce?
  \item \textbf{The hoist.} Identify the single line of the pipeline whose position
  (outside versus inside the terminal loop) is the structural signature of the
  fixed-operator corner. What goes wrong---in \emph{cost}, and separately in
  \emph{correctness of the abstraction}---if it is moved inside? (Contrast with the
  driven analysis of \cref{ch:driven}, where it belongs inside.)
  \item \textbf{Mirror, don't recompute.} A four-conductor problem needs the full
  $4\times4$ matrix $\mat{C}$. How many distinct bilinear evaluations
  $V_j^{\!\top}\Kop V_i$ does \code{gram\_reduce} actually perform, and which entries
  does \code{symmetric\_from\_upper} fill instead? What property of $\Kop$ licenses the
  mirroring?
  \item \textbf{Sign of the off-diagonal.} Using $q_i=C_{ii}V_i+\sum_{j\neq
  i}C_{ij}V_j$ and the physical picture of induced charge, argue that every
  off-diagonal entry of a Maxwell capacitance matrix is negative. Check your argument
  against the $\mat{C}$ of \cref{wex:es-twocond}.
  \item \textbf{The energy diagonal.} Redo \cref{wex:es-energy} for excitation~2:
  compute the field energy $U_e^{(2)}$ of $V_2$ directly and confirm it equals
  $\tfrac12 C_{22}$. Why must $C_{11}=C_{22}$ for this symmetric two-plate geometry?
  \item \textbf{Recover $C=\varepsilon A/d$.} In \cref{wex:es-twocond} the
  two-terminal capacitance came out $-C_{12}=\varepsilon A/L$. Explain why the
  physical between-plates capacitance is $-C_{12}$ (the negative of the mutual entry)
  rather than a diagonal entry, and relate this to the difference between the
  $\mat{C}$ matrix and the elementary single-number capacitor formula.
  \item \textbf{One stage swapped.} Using \cref{tab:es-swaps}, describe the
  magnetostatic analysis (\cref{ch:magnetostatics}) as ``the electrostatic analysis
  with these stages changed.'' Exactly which cells move, and which stay? Then do the
  same for the eigenmode analysis (\cref{ch:eigenmodes})---which cell is the
  most radical swap, and why?
  \item \textbf{The composition.} Write out the meaning of
  $\texttt{electrostatic}=\code{gram\_reduce}\circ\code{solve\_family}\circ
  \code{fe\_assemble}$ as data flowing left to right: name the type of the value
  passed from each stage to the next (config, operator, field family, matrix), and
  identify the one value---the operator $\Kop$---that is consumed by \emph{two}
  stages rather than one. Why does that shared consumption appear as the back-edge in
  \cref{fig:es-pipeline}?
\end{enumerate}
